\documentclass[a4paper,10pt]{article}
\usepackage[left=0.8cm, right=0.8cm, top=0.8cm, bottom=0.8cm]{geometry}
\usepackage{tikz}
\usepackage{amsmath,amssymb}
\usepackage{fontspec}
\usepackage{xcolor}
\usepackage{enumitem}
\usepackage{tabularx}
\usepackage{multicol}

% ── Fonts ──
\setmainfont{Helvetica Neue}[
  BoldFont = Helvetica Neue Bold,
  ItalicFont = Helvetica Neue Italic,
]

% ── Colors ──
\definecolor{case1bg}{HTML}{FFF9C4}    % yellow bg
\definecolor{case1border}{HTML}{F9A825} % yellow border
\definecolor{case1title}{HTML}{1565C0}  % blue title
\definecolor{case2bg}{HTML}{E8F5E9}    % green bg
\definecolor{case2border}{HTML}{2E7D32} % green border
\definecolor{case2title}{HTML}{2E7D32}  % green title
\definecolor{case3bg}{HTML}{E3F2FD}    % light blue bg
\definecolor{case3border}{HTML}{1565C0} % blue border
\definecolor{case3title}{HTML}{1565C0}  % blue title
\definecolor{case4bg}{HTML}{F3E5F5}    % purple bg
\definecolor{case4border}{HTML}{7B1FA2} % purple border
\definecolor{case4title}{HTML}{7B1FA2}  % purple title
\definecolor{formulabg}{HTML}{FFECB3}  % formula box bg
\definecolor{formulaborder}{HTML}{FF8F00}
\definecolor{examplebg}{HTML}{E8EAF6}
\definecolor{darktext}{HTML}{212121}
\definecolor{headerbg}{HTML}{1A237E}
\definecolor{accent}{HTML}{D32F2F}

\tikzset{
  casebox/.style={
    rounded corners=8pt,
    draw=#1,
    line width=1.5pt,
    inner sep=0pt,
  },
  casetitle/.style={
    font=\bfseries\normalsize,
    text=#1,
    align=center,
  },
  formulabox/.style={
    rounded corners=6pt,
    draw=formulaborder,
    fill=formulabg,
    line width=1pt,
    inner sep=6pt,
    align=center,
    font=\normalsize,
  },
  examplebox/.style={
    rounded corners=4pt,
    draw=case1border!60,
    fill=examplebg,
    line width=0.8pt,
    inner sep=5pt,
    align=left,
    font=\small,
  },
  % Box illustrations
  cardboard/.style={
    draw=brown!70!black,
    fill=brown!20,
    rounded corners=2pt,
    line width=0.8pt,
    minimum width=0.9cm,
    minimum height=0.7cm,
  },
  ball/.style={
    circle,
    draw=#1!70!black,
    fill=#1!40,
    minimum size=0.4cm,
    line width=0.6pt,
  },
}

\pagestyle{empty}

\begin{document}

% ╔══════════════════════════════════════════════════════════════╗
% ║                     MAIN TITLE                              ║
% ╚══════════════════════════════════════════════════════════════╝
\begin{tikzpicture}[remember picture, overlay]
  \fill[headerbg, rounded corners=8pt] 
    ([xshift=0.8cm, yshift=-0.4cm]current page.north west) 
    rectangle 
    ([xshift=-0.8cm, yshift=-1.7cm]current page.north east);
  \node[white, font=\LARGE\bfseries, anchor=center] 
    at ([yshift=-1.05cm]current page.north) 
    {DISTRIBUTING OBJECTS INTO BOXES};
\end{tikzpicture}

\vspace{1.2cm}

% ╔══════════════════════════════════════════════════════════════╗
% ║  2×2 GRID OF CASES                                          ║
% ╚══════════════════════════════════════════════════════════════╝

\noindent
\begin{tikzpicture}

% ─── Dimensions ───
\def\boxW{9.5cm}
\def\boxH{7.8cm}
\def\gap{0.4cm}

% ══════════════════════════════════════════════════
% CASE 1: Distinguishable Objects, Distinguishable Boxes (top-left)
% ══════════════════════════════════════════════════
\begin{scope}[shift={(0,0)}]
  \node[casebox=case1border, fill=case1bg, 
        minimum width=\boxW, minimum height=\boxH, anchor=north west] 
    (c1) at (0,0) {};
  
  % Title
  \node[casetitle=case1title, anchor=north, text width=8.5cm] 
    at ([yshift=-0.3cm]c1.north) 
    {1. DISTINGUISHABLE OBJECTS\\[-2pt] \& DISTINGUISHABLE BOXES};
  
  % Description
  \node[anchor=north west, text width=4.5cm, font=\small, text=darktext] 
    at ([xshift=0.4cm, yshift=-1.4cm]c1.north west) {
    Objects are \textbf{different}.\\
    Boxes are \textbf{different} (labeled).
  };
  
  % Balls
  \node[ball=red] at (1.0, -2.7) {};
  \node[ball=blue] at (1.7, -2.7) {};
  \node[ball=green!70!black] at (2.4, -2.7) {};
  \node[ball=orange] at (3.1, -2.7) {};
  \node[font=\scriptsize\bfseries, below] at (1.0,-2.95) {A};
  \node[font=\scriptsize\bfseries, below] at (1.7,-2.95) {B};
  \node[font=\scriptsize\bfseries, below] at (2.4,-2.95) {C};
  \node[font=\scriptsize\bfseries, below] at (3.1,-2.95) {D};
  
  % Arrows
  \draw[->, gray, thin] (1.0,-3.1) -- (1.2,-3.7);
  \draw[->, gray, thin] (1.7,-3.1) -- (2.4,-3.7);
  \draw[->, gray, thin] (2.4,-3.1) -- (2.4,-3.7);
  \draw[->, gray, thin] (3.1,-3.1) -- (3.6,-3.7);
  
  % Boxes
  \node[cardboard] at (1.2, -4.2) {};
  \node[font=\scriptsize, below] at (1.2,-4.65) {Box 1};
  \node[cardboard] at (2.4, -4.2) {};
  \node[font=\scriptsize, below] at (2.4,-4.65) {Box 2};
  \node[cardboard] at (3.6, -4.2) {};
  \node[font=\scriptsize, below] at (3.6,-4.65) {Box 3};
  
  % Interpretation
  \node[anchor=north west, text width=3.8cm, font=\small\itshape, text=case1title] 
    at (5.2, -1.4) {
    \textbf{Interpretation:}\\
    Each object can go\\to any box.
  };
  
  % Formula box
  \node[formulabox, anchor=center] at (6.8, -3.8) {
    \textbf{Total ways:}\\[3pt]
    $\displaystyle n^m$
  };
  
  % Legend
  \node[anchor=north west, font=\scriptsize, text=darktext] at (5.2, -5.1) {
    $n$ = number of boxes
  };
  \node[anchor=north west, font=\scriptsize, text=darktext] at (5.2, -5.45) {
    $m$ = number of objects
  };
  
  % Example
  \node[examplebox, anchor=south west, text width=8.3cm] 
    at ([xshift=0.4cm, yshift=0.25cm]c1.south west) {
    \textbf{Example:} $3^4 = 81$ ways
  };
\end{scope}

% ══════════════════════════════════════════════════
% CASE 2: Indistinguishable Objects, Distinguishable Boxes (top-right)
% ══════════════════════════════════════════════════
\begin{scope}[shift={(\boxW + \gap, 0)}]
  \node[casebox=case2border, fill=case2bg, 
        minimum width=\boxW, minimum height=\boxH, anchor=north west] 
    (c2) at (0,0) {};
  
  % Title
  \node[casetitle=case2title, anchor=north, text width=8.5cm] 
    at ([yshift=-0.3cm]c2.north) 
    {2. INDISTINGUISHABLE OBJECTS\\\& DISTINGUISHABLE BOXES};
  
  % Description
  \node[anchor=north west, text width=4.5cm, font=\small, text=darktext] 
    at ([xshift=0.4cm, yshift=-1.4cm]c2.north west) {
    Objects are \textbf{identical}.\\
    Boxes are \textbf{different} (labeled).
  };
  
  % Same-looking balls
  \foreach \x in {1.0, 1.7, 2.4, 3.1} {
    \node[ball=red] at (\x, -2.7) {};
  }
  
  % Stars and bars visual
  \draw[thick, darktext] (0.7,-3.4) -- (4.2,-3.4);
  \node[ball=red, minimum size=0.28cm] at (1.0, -3.4) {};
  \node[ball=red, minimum size=0.28cm] at (1.4, -3.4) {};
  \draw[very thick, case2border] (1.8,-3.1) -- (1.8,-3.7);
  \node[ball=red, minimum size=0.28cm] at (2.3, -3.4) {};
  \draw[very thick, case2border] (2.8,-3.1) -- (2.8,-3.7);
  \node[ball=red, minimum size=0.28cm] at (3.3, -3.4) {};
  \node[font=\tiny, text=case2title, anchor=west] at (3.6,-3.4) {$\leftarrow$ stars \& bars};
  
  % Boxes
  \node[cardboard] at (1.2, -4.2) {};
  \node[font=\scriptsize, below] at (1.2,-4.65) {Box 1};
  \node[cardboard] at (2.4, -4.2) {};
  \node[font=\scriptsize, below] at (2.4,-4.65) {Box 2};
  \node[cardboard] at (3.6, -4.2) {};
  \node[font=\scriptsize, below] at (3.6,-4.65) {Box 3};
  
  % Interpretation
  \node[anchor=north west, text width=4cm, font=\small\itshape, text=case2title] 
    at (5.0, -1.4) {
    \textbf{Interpretation:}\\
    Non-negative solutions of\\
    $x_1 + x_2 + \cdots + x_n = m$
  };
  
  % Formula box
  \node[formulabox, anchor=center, fill=green!10, draw=case2border] at (6.8, -3.8) {
    \textbf{Total ways:}\\[3pt]
    $\displaystyle \binom{m+n-1}{n-1}$
  };
  
  % Legend
  \node[anchor=north west, font=\scriptsize, text=darktext] at (5.0, -5.1) {
    $n$ = number of boxes
  };
  \node[anchor=north west, font=\scriptsize, text=darktext] at (5.0, -5.45) {
    $m$ = number of objects
  };
  
  % Example
  \node[examplebox, anchor=south west, text width=8.3cm, fill=green!5] 
    at ([xshift=0.4cm, yshift=0.25cm]c2.south west) {
    \textbf{Example:} $\binom{4+3-1}{3-1} = \binom{6}{2} = 15$ ways
  };
\end{scope}

% ══════════════════════════════════════════════════
% CASE 3: Distinguishable Objects, Indistinguishable Boxes (bottom-left)
% ══════════════════════════════════════════════════
\begin{scope}[shift={(0, -\boxH - \gap)}]
  \node[casebox=case3border, fill=case3bg, 
        minimum width=\boxW, minimum height=\boxH, anchor=north west] 
    (c3) at (0,0) {};
  
  % Title
  \node[casetitle=case3title, anchor=north, text width=8.5cm] 
    at ([yshift=-0.3cm]c3.north) 
    {3. DISTINGUISHABLE OBJECTS\\\& INDISTINGUISHABLE BOXES};
  
  % Description
  \node[anchor=north west, text width=4.5cm, font=\small, text=darktext] 
    at ([xshift=0.4cm, yshift=-1.4cm]c3.north west) {
    Objects are \textbf{different}.\\
    Boxes are \textbf{identical} (unlabeled).
  };
  
  % Different balls
  \node[ball=red] at (1.0, -2.7) {};
  \node[ball=blue] at (1.7, -2.7) {};
  \node[ball=green!70!black] at (2.4, -2.7) {};
  \node[ball=orange] at (3.1, -2.7) {};
  \node[font=\scriptsize\bfseries, below] at (1.0,-2.95) {A};
  \node[font=\scriptsize\bfseries, below] at (1.7,-2.95) {B};
  \node[font=\scriptsize\bfseries, below] at (2.4,-2.95) {C};
  \node[font=\scriptsize\bfseries, below] at (3.1,-2.95) {D};
  
  % Single unlabeled box
  \node[cardboard, minimum width=2.2cm, minimum height=0.9cm] at (2.0, -4.0) {};
  \node[font=\scriptsize, text=gray] at (2.0,-4.0) {\textit{groups}};
  
  % Interpretation
  \node[anchor=north west, text width=3.8cm, font=\small\itshape, text=case3title] 
    at (5.2, -1.4) {
    \textbf{Interpretation:}\\
    Partitioning objects\\into groups.
  };
  
  % Formula box
  \node[formulabox, anchor=center, fill=blue!5, draw=case3border] at (6.8, -3.5) {
    \textbf{Total ways:}\\[3pt]
    \textbf{Bell number}\\[1pt]
    $\displaystyle B_m$
  };
  
  % Stirling number note
  \node[anchor=north west, font=\scriptsize, text=darktext, text width=3.5cm] at (5.0, -4.8) {
    Into exactly $k$ groups:\\
    Stirling number $S(m,k)$\\[1pt]
    $B_m = \sum_{k=0}^{m} S(m,k)$
  };
  
  % Example
  \node[examplebox, anchor=south west, text width=8.3cm, fill=blue!3] 
    at ([xshift=0.4cm, yshift=0.25cm]c3.south west) {
    \textbf{Example:} $m=4$: $B_4 = 15$ ways
    {\scriptsize (Bell partitions of 4 objects)}
  };
\end{scope}

% ══════════════════════════════════════════════════
% CASE 4: Indistinguishable Objects, Indistinguishable Boxes (bottom-right)
% ══════════════════════════════════════════════════
\begin{scope}[shift={(\boxW + \gap, -\boxH - \gap)}]
  \node[casebox=case4border, fill=case4bg, 
        minimum width=\boxW, minimum height=\boxH, anchor=north west] 
    (c4) at (0,0) {};
  
  % Title
  \node[casetitle=case4title, anchor=north, text width=8.5cm] 
    at ([yshift=-0.3cm]c4.north) 
    {4. INDISTINGUISHABLE OBJECTS\\\& INDISTINGUISHABLE BOXES};
  
  % Description
  \node[anchor=north west, text width=4.5cm, font=\small, text=darktext] 
    at ([xshift=0.4cm, yshift=-1.4cm]c4.north west) {
    Objects are \textbf{identical}.\\
    Boxes are \textbf{identical} (unlabeled).
  };
  
  % Same balls
  \foreach \x in {1.0, 1.7, 2.4, 3.1} {
    \node[ball=red] at (\x, -2.7) {};
  }
  
  % Partitions of 4
  \node[anchor=north west, font=\small, text=darktext, text width=3.5cm] at (0.5, -3.4) {
    \textbf{Partitions of 4:}\\[2pt]
    {\small $4$\\
    $3 + 1$\\
    $2 + 2$\\
    $2 + 1 + 1$\\
    $1 + 1 + 1 + 1$}
  };
  
  % Interpretation
  \node[anchor=north west, text width=4cm, font=\small\itshape, text=case4title] 
    at (5.0, -1.4) {
    \textbf{Interpretation:}\\
    All objects and boxes\\are the same. Only\\count per group matters.
  };
  
  % Formula box
  \node[formulabox, anchor=center, fill=purple!5, draw=case4border] at (6.8, -3.5) {
    \textbf{Total ways:}\\[3pt]
    \textbf{Partitions of $m$}\\[1pt]
    $\displaystyle p(m)$
  };
  
  \node[anchor=north west, font=\scriptsize, text=darktext, text width=3.8cm] at (5.0, -4.8) {
    $p(m)$ = ways to write $m$ as\\
    sum of positive integers\\(order doesn't matter)
  };
  
  % Example
  \node[examplebox, anchor=south west, text width=8.3cm, fill=purple!3] 
    at ([xshift=0.4cm, yshift=0.25cm]c4.south west) {
    \textbf{Example:} $m=4$: $p(4) = 5$ ways
  };
\end{scope}

% ══════════════════════════════════════════════════
% QUICK REFERENCE TABLE (bottom)
% ══════════════════════════════════════════════════
\node[anchor=north, text width=18.5cm, align=center] 
  at (0.5*\boxW + 0.5*\gap, -2*\boxH - 2*\gap - 0.2cm) {
  \begin{tikzpicture}
    \node[rounded corners=6pt, draw=headerbg, fill=headerbg!6, line width=1.2pt, 
          inner sep=8pt, text width=18cm, align=center] {
      \textbf{\color{headerbg}\large Quick Reference — GATE DA}\\[6pt]
      \renewcommand{\arraystretch}{1.5}
      \begin{tabular}{|>{\bfseries}l|c|c|}
        \hline
        \rowcolor{headerbg!12}
        & \textbf{Distinguishable Boxes} & \textbf{Indistinguishable Boxes} \\
        \hline
        Distinguishable Objects & $n^m$ & Bell number $B_m$ \\
        \hline
        Indistinguishable Objects & $\displaystyle\binom{m+n-1}{n-1}$ {\small(Stars \& Bars)} & Partitions $p(m)$ \\
        \hline
      \end{tabular}\\[4pt]
      {\small $m$ = objects, $n$ = boxes \quad$\vert$\quad 
       \textcolor{accent}{\textbf{Most common in GATE:} Cases 1 \& 2}}
    };
  \end{tikzpicture}
};

\end{tikzpicture}

\end{document}
